% ============================================================
\chapter{REST API}
\label{app:api}
% ============================================================

Table~\ref{tab:api} lists the endpoints exposed by the NetMind backend. The
destructive actions reachable through the chat or the power endpoint produce
pending actions; only \code{/api/confirm} executes them, and it is called by
the interface, never by the language model. The server listens on
\code{127.0.0.1} and refuses requests from other origins.

{\small
\begin{xltabular}{\textwidth}{|L{5.4cm}|Y|}
\hline
\head \textbf{Method and path} & \textbf{Purpose} \\ \hline
\endhead
\code{POST /api/chat} & Run the agent loop for one user message; returns the reply, mode, evidence steps, and any pending actions \\ \hline
\code{POST /api/confirm} & Confirm or cancel a pending action (called by the interface only) \\ \hline
\code{GET /api/pending} & List pending actions \\ \hline
\code{GET /api/audit} & Audit log entries \\ \hline
\code{GET /api/devices} & Simulated devices with id, type, IP address, branch and status \\ \hline
\code{GET /api/alerts} & Alert history, newest first: devices going offline, degraded and back online (with downtime), firewall and camera findings \\ \hline
\code{POST /api/alerts/\{id\}/ack} & Acknowledge an alert (records that an operator has seen it; it stays open until the device recovers) \\ \hline
\code{GET /api/health} & The background collector's status and each device's confirmed state (up, degraded or down) \\ \hline
\code{GET /api/events} & Server-sent event stream telling the dashboard when alerts or device states change \\ \hline
\code{GET /api/logs} & Activity log: requested and confirmed actions, who made them, and their outcome \\ \hline
\code{GET /api/report} & Security and reliability summary \\ \hline
\code{POST /api/devices/\{id\}/power} & Turn a device on (runs) or off (creates a pending action) \\ \hline
\code{POST /api/notify-test} & Send a Telegram test message, if Telegram is configured \\ \hline
\code{GET /api/real-devices} & Real-network devices with IPv4, IPv6, MAC, hostname, status and \code{found\_by} (read-only) \\ \hline
\code{GET /api/real-alerts}, \code{/api/real-logs} & Real Network mode counterparts: alerts, and devices joining or leaving the network \\ \hline
\caption{REST API endpoints}\label{tab:api}\\
\end{xltabular}}

% ============================================================
\chapter{Running NetMind}
% ============================================================

From the repository root, inside WSL2 Ubuntu with a native docker-ce engine
and Containerlab installed:

\begin{enumerate}
\item Create a virtual environment and install the dependencies with
  \code{pip install -r requirements.txt}. Add \code{GROQ\_API\_KEY} to a
  \code{.env} file, and optionally \code{TELEGRAM\_BOT\_TOKEN} and
  \code{TELEGRAM\_CHAT\_ID} for notifications.
\item \code{make setup-perms} (once) gives the server the raw-socket
  capability and the re-attach helper, so it does not need to run as root.
\item \code{make up} builds the images, creates the Linux bridges, deploys the
  lab, and writes \code{registry.json}.
\item \code{make run} starts the web server, reachable at
  \code{http://127.0.0.1:8000}.
\item \code{make test} runs the automated tests, and \code{make lint} runs ruff.
\item \code{make down} removes the lab.
\end{enumerate}

% ============================================================
\chapter{Design Details}
\label{app:design}
% ============================================================

This appendix gives the detailed reasoning behind two parts of the design
that Chapter~5 summarises.

\section{Real-Network Discovery: Design Decisions}
\label{app:scan}

The decisions below are summarised in Table~\ref{tab:decisions}.

\textbf{Subnet detection.} The active interface is found by opening a UDP
socket towards a public address and reading the local address the operating
system assigns to it. No packet is actually sent, and the method relies only
on the Python standard library. If there is no default route at all, the
module reads the interface configuration with \code{ip addr} instead.
Alternatives were the netifaces and psutil libraries, which are more accurate
on machines with several interfaces but add an external dependency, and
gethostbyname applied to the machine's own hostname, which frequently returns
the loopback address 127.0.0.1 and is therefore unreliable. The chosen method
assumes a /24 subnet, which is a known limitation.

\textbf{Host discovery with ARP, Scapy and a ping sweep.} ARP was chosen
because it works at the data link layer: a device cannot ignore an ARP
request and still take part in an Ethernet network, so ARP finds devices that
are configured to block ICMP ping, and it also returns the MAC address. Scapy
was chosen to send the requests because it is a pure Python library that
integrates directly with the rest of the code and gives full control over
packet construction without an external program. In practice, however, some
devices answer ping but miss the short ARP window, so the ping sweep is not
only a fallback: running both and merging the results finds more devices than
either method alone, and \code{found\_by} makes visible which method saw each
device. The alternatives were Nmap, used directly or through python-nmap,
which is far more capable but is a separate binary larger than this module
requires; the arp-scan command-line tool, which is fast but is again an
external dependency; the operating system's ARP cache, which only lists
devices that the machine has already communicated with; and passive sniffing
of network traffic, which needs no probes but can take a long time to reveal
quiet devices.

\textbf{IPv6.} An IPv6 subnet is far too large to sweep, so the module sends
one ping to the IPv6 all-nodes multicast address on the active interface. Every
IPv6 device on the link answers, which places it in the kernel's neighbour
table; the module then reads that table and attaches each IPv6 address to the
device with the same MAC address. A device often has several IPv6 addresses
(a link-local \code{fe80::} address and one or more global ones), and all are
listed.

\textbf{Concurrency.} Checking 254 addresses one after another would take
several minutes, since each unanswered ping waits up to one second. A
ThreadPoolExecutor with 64 workers runs the pings in parallel and reduces the
sweep to a few seconds, and a second small pool runs the four discovery
methods at the same time, so a full scan takes about as long as its slowest
method. Threads suit this task because it is bound by waiting on the network
rather than by computation. The alternatives were multiprocessing, whose
separate processes add memory and start-up cost for work that is only
waiting, and asyncio, which would be efficient but would require the
surrounding code to be restructured around asynchronous functions.

\textbf{Hostname discovery with mDNS.} Multicast DNS was chosen because small
and home networks often have no reverse DNS records, so a reverse lookup
alone returns no name for most devices. mDNS lets printers, AirPlay devices,
Chromecast devices, and shared workstations announce their names without any
credentials or access to the router, and the zeroconf library provides a pure
Python implementation. The browse runs on every scan rather than once at
server start, so that a device that joins the network later is still found.
The alternatives were NetBIOS name queries, which are mostly limited to
Windows hosts; SSDP and UPnP discovery, which cover fewer device types; DHCP
lease tables read from the router, which need administrative access to the
router; and SNMP queries, which need SNMP to be enabled and credentials to be
known. A limitation of mDNS is that a device which advertises no service, such
as an idle phone, appears without a hostname, although it is still listed by
IP and MAC address.

\textbf{Reverse DNS and status checking.} Reverse DNS through the standard
library is used as a secondary source, because it is already available in
Python at no cost and gives names on networks that do maintain DNS records.
The status of a single device is determined with one ping, which keeps it
consistent with the ping sweep and needs no extra dependency. A TCP
connection check or an ARP probe would be more reliable against devices that
block ICMP, and are identified as possible improvements.

\textbf{Optional dependencies and output format.} Scapy and zeroconf are
imported inside try and except blocks, and their availability is recorded in
flags. If either library is missing the module degrades gracefully instead of
failing to start, which makes it easier to run on different machines. Results
are returned as plain dictionaries, rather than custom classes, because they
map directly to JSON for the API.

\section{Dashboard Refresh Performance}
\label{app:dash}

The first version of the dashboard polled the Docker daemon device by device,
which made one refresh take more than fifteen seconds with the full lab
running. Four changes made refreshes fast enough for regular polling:

\begin{enumerate}
\item all container states are read from a single container listing;
  per-container statistics, the slowest call and the only source of the CPU
  figures, are no longer requested, and the dashboard shows status rather than
  CPU usage;
\item the computed result is cached for three seconds, and concurrent callers
  share one computation instead of each starting their own;
\item firewall audits and camera credential checks used for the alert list
  are cached for ten seconds, and the cache is cleared whenever a rule,
  password or power state changes;
\item the frontend coalesces overlapping refreshes and releases a pressed
  button when its action finishes, so a button can no longer remain stuck.
\end{enumerate}
